% !PRINTED 32
\weberpage{32}

\begin{center}
{\large \S.\ 5.}

\vspace{0.5em}
{\large\gesperrt{Gebrochene Functionen; Theilbarkeit.}}
\end{center}

Wenn $F(x)$ und $f(x)$ zwei ganze rationale Functionen von $x$ sind,
so heisst der Bruch:
\begin{equation}
 {F(x)\over f(x)}
\tag{1}
\end{equation}
eine \gesperrt{gebrochene rationale} oder auch kurz
\gesperrt{gebrochene} oder \gesperrt{rationale Function} von $x$.

Ist der Grad des Zählers niedriger als der Grad des Nenners, so heisst
die Function \gesperrt{echt gebrochen}, im entgegengesetzten Falle
\gesperrt{unecht gebrochen}.

Nach \S~3 lassen sich die ganzen rationalen Functionen $Q$ und
$\varphi(x)$ so bestimmen, dass
\begin{equation}
 F(x)=Qf(x)+\varphi(x),
\tag{2}
\end{equation}
und der Grad von $\varphi(x)$ kleiner als der Grad von $f(x)$ ist;
demnach ist
\begin{equation}
 {F(x)\over f(x)}=Q+{\varphi(x)\over f(x)}
\tag{3}
\end{equation}
und daraus der Satz:

\gesperrt{Jede gebrochene Function kann in die Summe aus einer ganzen
und einer echt gebrochenen Function zerlegt werden.}

Ist $n$ der Grad von $f(x)$, so ist der Grad von $\varphi(x)$ höchstens
$n-1$, er kann aber auch niedriger sein; insbesondere kann auch der
Fall eintreten, dass $\varphi(x)$ \gesperrt{identisch verschwindet}.

\gesperrt{In diesem Falle heisst die Function $F(x)$ durch $f(x)$
theilbar.} Die Function
\[
 {F(x)\over f(x)}
\]
ist in diesem Falle nur scheinbar gebrochen, in Wirklichkeit der ganzen
Function $Q$ gleich.

Die Formel
\[
 {x^m-1\over x-1}=x^{m-1}+x^{m-2}+x^{m-3}+\cdots+1
\]
giebt hierfür ein einfaches Beispiel.
